\documentclass[aspectratio=1610]{beamer}
\usepackage[utf8]{inputenc}
\usepackage{sansmathfonts}
\usepackage[T1]{fontenc}

\usepackage[natbib=true,style=numeric,sorting=none]{biblatex}
\addbibresource{references.bib}

\usepackage{caption}
\usepackage{subcaption}
\usepackage{algorithm}
\usepackage{algpseudocode}
\usepackage{amsmath,amssymb,amsfonts}

\usepackage{tikz}
\usetikzlibrary{arrows.meta}
\usetikzlibrary{tikzmark}
\usetikzlibrary{backgrounds,calc,positioning}
\usetikzlibrary{shapes,decorations.pathreplacing}

\usepackage{listings}

\usetheme{main}

\captionsetup{font=scriptsize,labelfont=scriptsize}

\title[CM2 - Représentation d'informations]{CM2 - Représentation d'informations dans un ordinateur}

\date[]{}
\author[AM]{\textbf{Ammar Mian} \\ \footnotesize Associate professor, LISTIC, Université Savoie Mont Blanc \\ POLYTECH Annecy-Chambéry}

\newcommand{\red}[1]{\textcolor{red}{#1}}

%%%%% BASIC MATH DEFINITIONS %%%%%

\usepackage{amsmath,amsfonts,bm}

% Figure references
\def\figref#1{figure~\ref{#1}}
\def\Figref#1{Figure~\ref{#1}}

% Section references
\def\secref#1{section~\ref{#1}}
\def\Secref#1{Section~\ref{#1}}

% Algorithm references
\def\algref#1{algorithm~\ref{#1}}
\def\Algref#1{Algorithm~\ref{#1}}

% Common mathematical sets
\def\R{{\mathbb{R}}}
\def\N{{\mathbb{N}}}
\def\Z{{\mathbb{Z}}}
\def\C{{\mathbb{C}}}

% Common operations
\DeclareMathOperator{\Tr}{Tr}
\newcommand{\norm}[1]{\left\| #1 \right\|}
\newcommand{\argmin}[1]{\underset{#1}{\operatorname{argmin}}}
\newcommand{\argmax}[1]{\underset{#1}{\operatorname{argmax}}}

% Highlight new terms
\newcommand{\newterm}[1]{{\bf #1}}


\AtBeginBibliography{\scriptsize}

%%%%%%%%%%%%%%%%%%%%%%%%%%%%%%%%%%%%%%%%%%%%%%%%%%%%%%%%%%%%
%                     END OF PREAMBLE
%%%%%%%%%%%%%%%%%%%%%%%%%%%%%%%%%%%%%%%%%%%%%%%%%%%%%%%%%%%%

\begin{document}

%%%%%%%%%%%%%%%
% TITLE PAGE
%%%%%%%%%%%%%%%
\begin{frame}[noframenumbering,plain]
\titlepage
\end{frame}

\begingroup
\setbeamercolor{background canvas}{bg=main}
\setbeamercolor{titlelike}{fg=text-light, bg=light}
\begin{frame}[noframenumbering,plain]{Plan du cours}
    \tableofcontents[]
\end{frame}
\endgroup

\AtBeginSection[]{
    \setbeamercolor{background canvas}{bg=main}
    \begin{frame}[plain, noframenumbering]
        \tableofcontents[currentsection]
    \end{frame}
    \setbeamercolor{background canvas}{bg=light}
}

%%%%%%%%%%%%%%%%%%%%%%%%%%%%%%%%%%%%%%%%%%%%%%%%%%%%%%%%%%%%
%                     CONTENT STARTS HERE
%%%%%%%%%%%%%%%%%%%%%%%%%%%%%%%%%%%%%%%%%%%%%%%%%%%%%%%%%%%%

%%%%%%%%%%%%%%%
% SLIDE 2: Objectifs
%%%%%%%%%%%%%%%
\begin{frame}{Objectifs de ce cours}

  À l'issue de ce CM, vous serez capable de :

  \begin{itemize}
    \item Comprendre le principe du codage binaire
    \item Convertir des nombres entre différentes bases (2, 8, 10, 16)
    \item Représenter des entiers signés avec le complément à 2
    \item Comprendre la représentation des nombres flottants (IEEE 754)
    \item Connaître les principaux codages de caractères (ASCII, Unicode)
    \item Identifier les limites et pièges de ces représentations
  \end{itemize}

  \vspace{1em}

  \begin{alertblock}{Important}
    Ce cours nécessite un travail personnel pour approfondir les notions.
  \end{alertblock}

\end{frame}

%%%%%%%%%%%%%%%%%%%%%%%%%%%%%%%%%%%%%%%%%%%%%%%%%%%%%%%%%%%%
\section{Notion de Codage}
%%%%%%%%%%%%%%%%%%%%%%%%%%%%%%%%%%%%%%%%%%%%%%%%%%%%%%%%%%%%

%%%%%%%%%%%%%%%
% SLIDE 3: Qu'est-ce qu'un codage
egin{document}egin{frame}{Test}nd{frame}
\end{document}
